\section{Tóm tắt lý thuyết, giải pháp, thuật toán}

\subsection{Các lý thuyết và giải pháp liên quan}

Trong quá trình thực tập, em đã trực tiếp vận dụng các nền tảng của khoa học máy tính và nguyên lý thiết kế phần mềm vào quá trình phân tích, xây dựng và cải tiến hệ thống. Những kiến thức này được gắn với các bài toán thực tế và các giải pháp kỹ thuật được trình bày trong những phần tiếp theo.

\begin{itemize}[leftmargin=1.2cm, itemsep=3pt, topsep=3pt]
\item \textbf{Kiến trúc phân tầng và nguyên tắc đảo ngược phụ thuộc:} Hệ thống được tổ chức thành các tầng với trách nhiệm rõ ràng theo nguyên lý kiến trúc sạch (Clean Architecture) [1]: tầng giao diện xử lý tương tác với người dùng, tầng ứng dụng và nghiệp vụ xử lý các quy tắc của hệ thống, còn tầng hạ tầng đảm nhiệm việc làm việc với cơ sở dữ liệu và các thành phần bên ngoài. Các tầng nghiệp vụ cốt lõi được thiết kế để không phụ thuộc trực tiếp vào giao diện hay công nghệ lưu trữ cụ thể. Nhờ đó, các quy tắc nghiệp vụ quan trọng được tập trung và quản lý thống nhất, đồng thời có thể được sử dụng lại ở nhiều chức năng khác nhau mà không phải sao chép hoặc viết lại cùng một logic.

\item \textbf{Phân tách lệnh và truy vấn (CQRS):} Hệ thống áp dụng mô hình phân tách lệnh và truy vấn (CQRS) [4] để tách riêng các thao tác làm thay đổi dữ liệu (Command) với các thao tác chỉ đọc và truy xuất dữ liệu (Query). Với các nghiệp vụ như đặt phòng, thay vì thay đổi trực tiếp trạng thái của bản ghi, hệ thống xử lý thông qua các lệnh thể hiện rõ mục đích của nghiệp vụ. Cách tiếp cận này giúp dễ dàng theo dõi và lưu lại lịch sử thay đổi, đồng thời cho phép phần xử lý dữ liệu đọc được tối ưu riêng cho các chức năng như báo cáo và thống kê mà không ảnh hưởng đến quá trình ghi dữ liệu.

\item \textbf{Chi phí của tầng trừu tượng khi truy cập dữ liệu (ORM Overhead):} Các công cụ ánh xạ đối tượng-quan hệ (ORM) như Entity Framework [6] giúp lập trình viên thao tác với cơ sở dữ liệu thông qua ngôn ngữ bậc cao. Tuy nhiên, nếu sử dụng không phù hợp, tầng trừu tượng này có thể sinh ra các câu truy vấn phức tạp, nạp thừa dữ liệu qua mạng (over-fetching) hoặc dẫn đến hiện tượng bùng nổ tích Descartes khi nối nhiều bảng quan hệ. Tại các tác vụ xử lý lượng bản ghi lớn, việc xem xét kế hoạch thực thi (Execution Plan) và cơ chế khóa [7] kết hợp sử dụng các câu truy vấn định hướng (Projection Queries) là yếu tố quan trọng để đảm bảo hiệu năng.

\item \textbf{Khóa phân tán và quản lý vòng đời khóa (Distributed Lock with TTL):} Khi nhiều tiến trình cùng truy cập và cập nhật tài nguyên dùng chung, khóa phân tán trên Redis hỗ trợ xử lý tuần tự và hạn chế xung đột tài nguyên [8]. Việc thiết lập thời gian sống (TTL) cho khóa là điều cần thiết để giảm nguy cơ khóa bị giữ vô thời hạn (stale lock) trong trường hợp một tiến trình gặp sự cố bất thường trước khi kịp giải phóng khóa.

\item \textbf{Cơ chế phát hiện thay đổi trên giao diện (Change Detection):} Trong Angular, quá trình phát hiện thay đổi có thể khiến các biểu thức trong template được đánh giá lại qua nhiều chu kỳ cập nhật giao diện [5]. Do đó, việc gọi trực tiếp các hàm tính toán phức tạp hoặc quét mảng trong template có thể bị kích hoạt lặp đi lặp lại, gây sụt giảm khung hình và làm giao diện phản hồi chậm chạp.

\item \textbf{Thuật toán gộp khoảng và tính giao thời gian (Interval Merging):} Cho một tập các khoảng thời gian có khả năng chồng lấn $I = \{[s_i, e_i]\}$ [3], thuật toán tiến hành sắp xếp theo mốc bắt đầu với độ phức tạp $O(k \log k)$, sau đó duyệt tuần tự trong $O(k)$ để hợp nhất các khoảng có phần giao nhau thành các khoảng liên tục không chồng lấn. Từ các khoảng đã gộp, hệ thống tính toán phần giao thực tế với một khung thời gian quan sát $[T_{\text{start}}, T_{\text{end}}]$ cho trước.

\item \textbf{Hàm băm một chiều bảo vệ dữ liệu định danh (One-way Hashing):} Nhằm tuân thủ các quy chuẩn bảo mật và an toàn thông tin cá nhân (PII), dữ liệu định danh của khách lưu trú được chuyển thành giá trị băm có độ dài cố định bằng hàm băm một chiều. Kỹ thuật này cho phép hệ thống đối chiếu và gom nhóm khách hàng dựa trên giá trị định danh mà không cần sử dụng trực tiếp thông tin gốc trên phân hệ báo cáo.

\item \textbf{Chuẩn hóa chuỗi văn bản tìm kiếm đa ngôn ngữ (Unicode Normalization):} Xử lý tìm kiếm trong môi trường đa ngôn ngữ đòi hỏi chuẩn hóa cả chuỗi người dùng nhập lẫn dữ liệu nguồn về dạng chuẩn (như NFC). Đồng thời, thuật toán sử dụng các thuộc tính ký tự Unicode (\texttt{\textbackslash p\{L\}}, \texttt{\textbackslash p\{N\}}) theo chuẩn Unicode [9] thay cho các biểu thức chính quy thuần ASCII, nhằm bảo toàn trọn vẹn dấu tiếng Việt và các hệ ký tự tượng hình như tiếng Hàn và tiếng Trung.
\end{itemize}

\subsection{Cách giải quyết của em}

\subsubsection{Tính năng khóa phòng: xử lý thay đổi trạng thái}

Tính năng khóa phòng phục vụ nhu cầu tạm ngừng kinh doanh một hoặc nhiều phòng trong một khoảng thời gian xác định kèm lý do bảo trì hoặc sửa chữa. Vấn đề phức tạp nhất xuất hiện ở thao tác chỉnh sửa: người dùng mở một lượt khóa hiện có, sau đó có thể đồng thời thêm phòng mới, loại bỏ phòng cũ, thay đổi khung thời gian bắt đầu/kết thúc và cập nhật lại nội dung lý do.

Phương án đơn giản nhất có thể nghĩ tới là xóa sạch toàn bộ lượt khóa cũ rồi tạo lại lượt mới từ đầu. Phương án này dễ lập trình nhưng bộc lộ hai nhược điểm: lịch sử biến động phòng bị xóa trắng khiến ban quản lý khó truy vết những thay đổi đã thực hiện; đồng thời, các phòng không hề bị thay đổi cũng phải trải qua chu kỳ hủy và tạo lại, gây ra các biến động trạng thái dư thừa và làm tăng nguy cơ xung đột khóa với các đơn đặt phòng khác đang diễn ra.

Để xử lý vấn đề này, em so sánh trạng thái phòng trước và sau khi chỉnh sửa để xác định chính xác những thay đổi thực tế cần xử lý. Cụ thể, gọi $R_{\text{old}}$ là tập phòng đang khóa và $R_{\text{new}}$ là tập phòng sau chỉnh sửa: các phòng thuộc $R_{\text{new}} \setminus R_{\text{old}}$ sinh lệnh giữ phòng mới; các phòng thuộc $R_{\text{old}} \setminus R_{\text{new}}$ sinh lệnh giải phóng; nếu khung thời gian thay đổi thì sinh lệnh điều chỉnh thời gian; và nếu chỉ sửa đổi lý do thì chỉ phát sinh lệnh cập nhật ghi chú. Thành phần xử lý sai khác này được thiết kế độc lập với cơ sở dữ liệu, giúp thuận tiện cho việc viết kiểm thử tự động với các kịch bản tổ hợp.

\begin{figure}[htbp]
\centering
\begin{tikzpicture}[
  hop/.style={draw=xamvien, fill=xamnhat, rounded corners=2pt, align=center,
              font=\small, inner sep=4pt},
  lenh/.style={draw=xanhkv, fill=white, rounded corners=2pt, align=left,
               font=\scriptsize, inner sep=4pt, text width=4.2cm},
  mui/.style={-{Stealth[length=2mm]}, draw=xamvien, thick}
]
\node[hop, text width=3.6cm, minimum height=1.1cm] (cu) at (0,0)
  {\textbf{Trạng thái hiện lưu}\\[1pt]{\scriptsize Danh sách phòng và khoảng thời gian}};
\node[hop, text width=3.6cm, minimum height=1.1cm] (moi) at (0,-2.2)
  {\textbf{Trạng thái gửi lên}\\[1pt]{\scriptsize Dữ liệu sau khi người dùng sửa}};
\node[hop, text width=2.6cm, minimum height=1.5cm, fill=white, draw=xanhkv, thick] (ss) at (5.0,-1.1)
  {\textbf{So sánh sai khác}\\[2pt]{\scriptsize Xác định biến động thực tế}};

\node[lenh] (l1) at (10.1,0.6) {Phòng mới bổ sung\\ $\rightarrow$ Lệnh giữ phòng};
\node[lenh] (l2) at (10.1,-0.75) {Phòng bị loại bỏ\\ $\rightarrow$ Lệnh giải phóng giữ phòng};
\node[lenh] (l3) at (10.1,-2.1) {Khung thời gian thay đổi\\ $\rightarrow$ Lệnh điều chỉnh thời gian};
\node[lenh] (l4) at (10.1,-3.45) {Chỉ lý do thay đổi\\ $\rightarrow$ Lệnh cập nhật ghi chú};

\draw[mui] (cu.east) -- (ss.west);
\draw[mui] (moi.east) -- (ss.west);
\draw[mui] (ss.east) -- (l1.west);
\draw[mui] (ss.east) -- (l2.west);
\draw[mui] (ss.east) -- (l3.west);
\draw[mui] (ss.east) -- (l4.west);
\end{tikzpicture}
\caption{Chỉ sinh lệnh cho phần thực sự thay đổi thay vì huỷ và tạo lại toàn bộ}
\label{hinh:sinhlenh}
\end{figure}

\indent Hình \ref{hinh:sinhlenh} minh họa trực quan luồng phát sinh lệnh có chọn lọc. Toàn bộ các quy tắc kiểm tra tính hợp lệ của dữ liệu đầu vào cũng được tập trung tại một đầu mối duy nhất: loại khóa hợp lệ, giới hạn độ dài lý do, bắt buộc chọn ít nhất một phòng, không cho phép chọn trùng lặp và thời điểm kết thúc luôn phải sau thời điểm bắt đầu. Việc gom cụm quy tắc này bảo đảm thông báo lỗi đồng nhất trên toàn hệ thống.

Ở tầng giao diện, em xây dựng hộp thoại thao tác tích hợp vào sơ đồ lịch phòng và bổ sung cơ chế kiểm tra chặn khi đặt phòng. Đặc biệt, sau khi hoàn thành tính năng, em chủ động phân tích tác động dây chuyền và phát hiện các công thức tính tỷ lệ công suất phòng đang tính cả phòng bị khóa vào số phòng khả dụng để bán. Điều này khiến tỷ lệ lấp đầy hiển thị thấp hơn thực tế, làm sai lệch đánh giá hiệu quả kinh doanh của cơ sở. Em đã chủ động trao đổi và điều chỉnh lại công thức để loại trừ các phòng đang khóa bảo trì ra khỏi mẫu số tính công suất, giúp số liệu thống kê phản ánh chính xác thực tế vận hành.

\subsubsection{Thống kê lượt khách nghỉ qua đêm và tổng số đêm nghỉ}

Yêu cầu bài toán là bổ sung hai chỉ tiêu thống kê quan trọng vào báo cáo phòng trong một khoảng thời gian tùy chọn. Khi phân tích yêu cầu, em nhận thấy phần tính toán cần xử lý ba vấn đề chính:

\textbf{1. Xử lý các khoảng lưu trú chồng lấn:} Một khách hàng có thể có nhiều bản ghi lưu trú nối tiếp hoặc chồng lấn nhau trong kỳ do đổi phòng hoặc gia hạn phòng. Nếu cộng dồn số đêm của từng bản ghi, các ngày trùng lặp sẽ bị tính nhiều lần làm phồng số liệu. Em áp dụng thuật toán gộp khoảng [3]: sắp xếp toàn bộ các bản ghi của cùng một khách theo thời điểm nhận phòng, sau đó duyệt tuần tự để hợp nhất các khoảng có phần giao nhau thành các khoảng lưu trú liên tục duy nhất.

\textbf{2. Cắt khoảng theo biên kỳ báo cáo:} Một lượt lưu trú có thể bắt đầu trước hoặc kết thúc sau khoảng thời gian đang xét. Em xác định riêng quy tắc tính toán cho từng chỉ tiêu: số lượt khách chỉ được ghi nhận khi thời điểm nhận phòng nằm trong kỳ báo cáo (vì lượt là một sự kiện phát sinh); trong khi tổng số đêm nghỉ được tính bằng độ dài phần giao thực tế giữa các khoảng lưu trú đã gộp với khung thời gian báo cáo.

\begin{figure}[htbp]
\centering
\begin{tikzpicture}[font=\small]
\draw[->, thick, draw=xamvien] (0,0) -- (11.9,0) node[right,font=\scriptsize]{thời gian};
\foreach \x/\l in {1/1,3/3,5/5,7/7,9/9,11/11} {
  \draw[draw=xamvien] (\x,-0.1) -- (\x,0.1);
  \node[font=\scriptsize, below=1pt] at (\x,-0.1) {ngày \l};
}
\fill[xanhkv, opacity=0.08] (3,-0.55) rectangle (9,4.3);
\draw[dashed, thick, draw=xanhkv] (3,-0.55) -- (3,4.3);
\draw[dashed, thick, draw=xanhkv] (9,-0.55) -- (9,4.3);
\node[font=\scriptsize\bfseries, text=xanhkv] at (6,4.05) {kỳ báo cáo đang xem};

\node[font=\scriptsize, anchor=east] at (-0.15,3.4) {bản ghi 1};
\draw[line width=3.5pt, draw=xamvien] (1,3.4) -- (4,3.4);
\node[font=\scriptsize, anchor=east] at (-0.15,2.85) {bản ghi 2};
\draw[line width=3.5pt, draw=xamvien] (3.4,2.85) -- (6,2.85);
\node[font=\scriptsize, anchor=east] at (-0.15,2.3) {bản ghi 3};
\draw[line width=3.5pt, draw=xamvien] (7,2.3) -- (11,2.3);

\node[font=\scriptsize, anchor=east] at (-0.15,1.45) {sau khi gộp};
\draw[line width=3.5pt, draw=xanhkv] (1,1.45) -- (6,1.45);
\draw[line width=3.5pt, draw=xanhkv] (7,1.45) -- (11,1.45);

\node[font=\scriptsize, anchor=east] at (-0.15,0.75) {phần được tính};
\draw[line width=3.5pt, draw=xanhkv] (3,0.75) -- (6,0.75);
\draw[line width=3.5pt, draw=xanhkv] (7,0.75) -- (9,0.75);
\end{tikzpicture}
\caption{Gộp các khoảng lưu trú chồng lấn rồi cắt theo biên của kỳ báo cáo}
\label{hinh:gopkhoang}
\end{figure}

\indent Hình \ref{hinh:gopkhoang} mô tả hai bước xử lý: ba bản ghi ban đầu sau khi hợp nhất tạo thành hai khoảng rời rạc; khi đối chiếu theo khung báo cáo, các phần nằm ngoài khoảng thời gian báo cáo được loại bỏ. Bản ghi 1 bắt đầu trước kỳ nên không tính vào số lượt, nhưng các đêm nằm trong kỳ vẫn được tính đầy đủ vào tổng số đêm nghỉ.

\textbf{3. Xử lý trường hợp định danh rỗng:} Khách hàng được nhận diện thông qua giá trị băm của thông tin định danh cá nhân. Trong quá trình kiểm tra thực tế, em phát hiện một vấn đề nghiệp vụ: đối với những khách vãng lai hoặc khách thuê theo giờ không khai báo giấy tờ tùy thân, chuỗi định danh ban đầu là chuỗi rỗng. Khi đi qua hàm băm, chuỗi rỗng luôn sinh ra cùng một giá trị băm cố định. Nếu gom nhóm đơn thuần theo giá trị băm này, toàn bộ khách không khai báo giấy tờ trong suốt kỳ báo cáo sẽ bị gộp thành một người duy nhất, dẫn đến số lượt khách hiển thị bị sai lệch nghiêm trọng. Để khắc phục, hệ thống phân biệt trường hợp không có định danh và quản lý các lượt lưu trú này dựa trên mã đặt phòng (booking code) kết hợp thời gian lưu trú cụ thể, nhờ đó mỗi lượt khách chưa khai báo giấy tờ vẫn được tính toán độc lập và chính xác theo đúng nghiệp vụ.

Toàn bộ logic tính toán được đóng gói thành một lớp thuần túy độc lập với cơ sở dữ liệu, cho phép thiết lập bộ kiểm thử bao phủ các tình huống biên phức tạp. Trong đó, em bổ sung trường hợp kiểm thử kiểm tra tính chất nhất quán của thuật toán: tổng kết quả tính toán của hai khoảng thời gian liên tiếp ghép lại phải bằng đúng kết quả khi tính một lần trên toàn bộ khoảng thời gian gộp.

\subsubsection{Tối ưu quá trình xuất dữ liệu và xử lý các vấn đề về hiệu năng}

Chức năng xuất tệp sổ quỹ gặp tình trạng chậm khi dữ liệu lớn. Phân tích mã nguồn cho thấy phần lấy tên phòng được xây dựng bằng nhiều biểu thức lồng nhau, thực hiện phép nối trên nhiều bảng và lọc trùng lặp qua nhiều bước. Tuy nhiên, điểm nghẽn lớn nhất nằm ở thuật toán xử lý dữ liệu sau khi nạp: mã nguồn cũ duyệt qua từng dòng của tệp và quét tuần tự toàn bộ danh sách phòng để tìm kiếm tên tương ứng. Với tệp $n$ dòng và danh sách phòng $m$ phần tử, số phép so sánh có độ phức tạp $O(n \times m)$, gây độ trễ đáng kể khi xuất hàng nghìn bản ghi.

Em tái cấu trúc thuật toán theo ba bước: kiểm tra sớm nếu không có dữ liệu cần xuất thì dừng ngay; chuyển đổi chuỗi biểu thức phức tạp thành một câu lệnh truy vấn chuẩn hóa duy nhất ở tầng cơ sở dữ liệu để lấy đủ thông tin cần thiết; và chuyển danh sách phòng thành một bảng tra cứu (Hash Map) theo khóa định danh. Thuật toán mới giảm độ phức tạp xuống $O(n + m)$, giảm số vòng lặp và làm luồng xử lý đơn giản hơn, giúp thời gian xử lý khi xuất dữ liệu lớn được cải thiện rõ rệt.

Một tình huống khác liên quan đến tranh chấp tài nguyên cơ sở dữ liệu khi người dùng xuất báo cáo hóa đơn chi tiết. Thay vì can thiệp vào cơ chế quản lý giao dịch chung vốn tiềm ẩn nhiều rủi ro, em phân tích câu truy vấn và nhận thấy hệ thống đang nạp kèm một số dữ liệu quan hệ dư thừa mà mẫu tệp không sử dụng. Việc nạp dữ liệu không cần thiết làm câu truy vấn phức tạp hơn, tăng lượng dữ liệu xử lý và kéo dài thời gian thực thi, từ đó làm tăng khả năng tranh chấp tài nguyên trên cơ sở dữ liệu [7]. Việc loại bỏ các quan hệ nạp thừa và tối ưu danh sách truy vấn đã giúp xử lý được hiện tượng nghẽn này mà chỉ cần điều chỉnh một phần nhỏ mã nguồn.

Ngoài ra, trong quá trình rà soát, em nhận thấy một số khóa phân tán trên Redis được tạo ra nhưng chưa thiết lập thời gian sống (TTL) [8]. Nếu tiến trình xử lý gặp sự cố bất thường trước khi kịp giải phóng, khóa có thể bị giữ lại và gây ảnh hưởng đến các tiến trình tiếp theo. Em đã bổ sung cơ chế TTL tự động giải phóng khóa, giúp giảm thiểu nguy cơ tài nguyên bị chiếm dụng vô thời hạn.

\subsubsection{Tập trung quy tắc phân quyền vào một chỗ}

Yêu cầu đặt ra là phân tách quyền hủy đơn phòng cũ thành hai quyền hạn độc lập dựa trên trạng thái lưu trú của khách. Bài toán có hai khía cạnh cần xử lý đồng bộ:

Về mặt dữ liệu, hàng nghìn khách hàng đang vận hành đã được gán quyền hạn theo cấu trúc cũ. Việc chuyển đổi sang hai quyền mới phải đảm bảo tính liên tục, không làm gián đoạn hay mất mát quyền của bất kỳ tài khoản nào. Em xây dựng kịch bản chuyển đổi dữ liệu tự động, đồng thời viết kèm kịch bản khôi phục (rollback) để đưa dữ liệu về trạng thái ban đầu nếu phát sinh sự cố. Kịch bản dự phòng là nguyên tắc bắt buộc khi can thiệp vào dữ liệu sản xuất.

Về mặt giao diện, logic kiểm tra quyền xuất hiện tại nhiều vị trí trên màn hình quầy. Thay vì rải các câu điều kiện kiểm tra ở từng nút bấm, em đóng gói toàn bộ quy tắc vào một dịch vụ quản lý quyền chuyên biệt. Các thành phần hiển thị chỉ cần gọi một phương thức duy nhất để xác định quyền thao tác của người dùng hiện tại. Đồng thời, quyền thao tác của người dùng phải được kiểm tra ở phía máy chủ, bởi việc ẩn hoặc vô hiệu hóa nút trên giao diện chỉ đóng vai trò hỗ trợ trải nghiệm chứ không thể thay thế cho việc kiểm tra quyền thực tế.

\subsubsection{Giữ cho giao diện nhẹ và tách logic khỏi phần hiển thị}

Quy ước của nhóm dự án là hạn chế tối đa việc gọi các hàm hoặc thuộc tính tính toán trực tiếp trong template của Angular [5]. Em áp dụng hai giải pháp: gán trước kết quả tính toán vào biến trạng thái khi dữ liệu thay đổi để giao diện chỉ việc đọc giá trị; và sử dụng các pipe thuần túy (pure pipes) có cơ chế lưu kết quả để tái sử dụng (memoization) [5] cho các bảng dữ liệu danh sách dài, giúp hệ thống chỉ tính toán lại khi tham số đầu vào thực sự biến động.

Tương tự, khi xây dựng ô tìm kiếm cho tài liệu hướng dẫn tích hợp kênh bán phòng trực tuyến, em tách riêng logic xử lý chuỗi ra một lớp tiện ích gồm ba bước: chuẩn hóa văn bản, tách mảng từ khóa và đối chiếu nội dung. Khi chuẩn hóa, em sử dụng bộ nhận diện ký tự Unicode thay vì lọc theo bảng chữ cái ASCII nhằm bảo toàn dấu tiếng Việt và chữ tượng hình [9]; khi đối chiếu, hệ thống yêu cầu tiêu đề chứa đầy đủ các từ khóa mà không bắt buộc theo đúng thứ tự gõ.

\subsubsection{Dựng lại biểu đồ doanh thu và một lỗi lặp lại ba lần}

Khi xây dựng lại màn hình quản lý tổng quan, em tiếp nhận thành phần biểu đồ doanh thu cũ vốn trộn lẫn logic nạp dữ liệu, tính toán thống kê và dựng đồ họa trong một tệp duy nhất. Em đã tái cấu trúc thành ba thành phần độc lập, chuyển phần giao tiếp dữ liệu về đúng tầng dịch vụ chuyên trách.

Trong quá trình hoàn thiện, ba phiếu lỗi riêng biệt xuất hiện trên cùng màn hình: bỏ chọn một chi nhánh nhưng tổng doanh thu không đổi, biểu đồ cơ cấu nguồn khách không vẽ lại, và khối tổng tiền phòng giữ nguyên giá trị cũ. Qua việc kiểm tra ba lỗi trên, em nhận thấy chúng đều bắt nguồn từ cùng một nguyên nhân: các thành phần con tự lưu trữ bản sao danh sách chi nhánh riêng biệt, dẫn đến sự lệch pha khi có sự kiện cập nhật. Em khắc phục bằng cách chuyển sang chia sẻ một nguồn dữ liệu trạng thái chi nhánh chung (Single Source of Truth). Sau khi thống nhất nguồn dữ liệu trạng thái ở tầng dịch vụ dùng chung, ba trường hợp lỗi trên không còn xuất hiện trong quá trình kiểm thử. Bài học rút ra là khi một hiện tượng xuất hiện ở nhiều nơi, việc tìm đúng nguyên nhân gốc trong thiết kế luồng dữ liệu sẽ giúp xử lý vấn đề hiệu quả hơn so với việc sửa riêng lẻ tại từng giao diện.

\subsection{Liên hệ và so sánh với các cách đã có}

Trong phạm vi công việc được giao, em chủ yếu so sánh giải pháp được lựa chọn với các phương án đã cân nhắc trong quá trình phát triển thực tế tại doanh nghiệp, thay vì so sánh với các sản phẩm bên ngoài.

\textbf{Về giải pháp tính năng khóa phòng:} Phương án ban đầu có thể nghĩ tới là xóa sạch lượt khóa cũ rồi tạo lại, ưu điểm là mã nguồn ngắn và triển khai nhanh. Tuy nhiên, giải pháp so sánh trạng thái trước và sau khi chỉnh sửa mà em áp dụng mang lại nhiều ưu thế hơn: lưu giữ được đầy đủ nhật ký thao tác cho người quản trị; giảm thiểu biến động trạng thái trên hệ thống; và hạn chế nguy cơ xung đột khóa dữ liệu với các đơn đặt phòng khác đang diễn ra.

\textbf{Về vị trí đặt phần tính toán thống kê:} Ban đầu, em đặt module tính toán tại dịch vụ quản lý do báo cáo được mở từ giao diện này. Sau khi trao đổi cùng cán bộ hướng dẫn, em nhận ra cách tiếp cận đó chưa hợp lý: dữ liệu lưu trú thuộc trách nhiệm quản lý của dịch vụ đặt phòng, việc để dịch vụ quản lý tự đọc và tính toán sẽ tạo ra hai nơi cùng định nghĩa một chỉ tiêu nghiệp vụ. Em đã chủ động hủy bỏ phần mã cũ để chuyển toàn bộ logic tính toán về dịch vụ đặt phòng. Nguyên tắc em rút ra là trong kiến trúc của hệ thống, nhóm ưu tiên đặt logic xử lý tại dịch vụ sở hữu dữ liệu nguồn.

\textbf{Về phạm vi tác động của một thay đổi:} Khi phát hiện lỗi định dạng ngày trên giao diện tiếng Hàn, em nhận thấy mã nguồn xử lý ngày tháng cũ bị phân tán và ban đầu đã viết lại hơn một trăm dòng mã trên nhiều tệp khác nhau. Tuy nhiên, nhận thấy việc thay đổi trên diện rộng tiềm ẩn rủi ro khi chưa có đầy đủ kiểm thử tự động bao phủ toàn bộ các màn hình, em đã hủy bỏ thay đổi đó và chuyển sang giải pháp an toàn hơn: chỉ tối ưu hàm đọc ngày tháng tại đúng điểm phát sinh lỗi trong phạm vi khoảng hai mươi dòng mã. Bài học em rút ra là việc tái cấu trúc mã nguồn trên diện rộng cần được hoạch định độc lập với bộ kiểm thử bảo vệ [2], không nên tùy tiện gộp chung vào một phiếu sửa lỗi nóng cần phát hành gấp.
